%
\usepackage[T1]{fontenc}
% T1 fonts will be used to generate the final print and online PDFs,
% so please use T1 fonts in your manuscript whenever possible.
% Other font encondings may result in incorrect characters.
%
\usepackage{graphicx}
% Used for displaying a sample figure. If possible, figure files should
% be included in EPS format.
%
% If you use the hyperref package, please uncomment the following two lines
% to display URLs in blue roman font according to Springer's eBook style:
%\usepackage{color}
%\renewcommand\UrlFont{\color{blue}\rmfamily}
%\urlstyle{rm}
%

\usepackage{comment}
\usepackage{tabularx}
\usepackage{todonotes} 

\usepackage{listings} 
\lstset{
  basicstyle=\footnotesize\tt,        % the size of the fonts that are used for the code
  breakatwhitespace=false,         % sets if automatic breaks should only happen at whitespace
  breaklines=true,                 % sets automatic line breaking
  captionpos=b,                    % sets the caption-position to bottom
  extendedchars=true,              % lets you use non-ASCII characters; for 8-bits encodings only, does not work with UTF-8
  frame=single,                    % adds a frame around the code
  language=Java,                 % the language of the code
  keywordstyle=\bf,
  showspaces=false,                % show spaces everywhere adding particular underscores; it overrides 'showstringspaces'
  showstringspaces=false,          % underline spaces within strings only
  showtabs=false,                  % show tabs within strings adding particular underscores
  tabsize=2                       % sets default tabsize to 2 spaces
}

% for subfigure
% warning ist ein bekanntes Problem
\usepackage{subcaption}

\usepackage{comment}


\newcommand{\patternname}[1]{\newpage{\large{\textsc{#1}}}}

\newcommand{\patternnameinline}[1]{\textsc{#1}}


%\textbf{\textsc{Foo}}

\newcommand{\patternOne}{
    \patternname{Pattern 1: Core of Content}
}

\newcommand{\patternTwo}{
    \patternname{Pattern 2: Accessibility of Diagrams}
}

\newcommand{\patternThree}{
    \patternname{Pattern 3: Accessibility of Graphics}
}

\newcommand{\patternFour}{
    \patternname{Pattern 4: Sonification of Time Series}
}

\newcommand{\patternFive}{
    \patternname{Pattern 5: Graphics with Generated Explanation}
}

\newcommand{\patternSix}{
    \patternname{Pattern 6: Accessibility-Aware AI Tutor}
}

\newcommand{\patternSeven}{
    % KI: C.10 benennt Pattern 7 um; der neue Name betont die strukturierte Begleitfunktion statt generischer Podcast-Erzeugung.
    \patternname{Pattern 7: Structured Audio Companion}

}

\newcommand{\patternOneInline}{
    \patternnameinline{Core of Content}
}

\newcommand{\patternTwoInline}{
    \patternnameinline{Accessibility of Diagrams}
}

\newcommand{\patternThreeInline}{
    \patternnameinline{Accessibility of Graphics}
}

\newcommand{\patternFourInline}{
    \patternnameinline{Sonification of Time Series}
}

\newcommand{\patternFiveInline}{
    \patternnameinline{Graphics with Generated Explanation}
}

\newcommand{\patternSixInline}{
    \patternnameinline{Accessibility-Aware AI Tutor}
}

\newcommand{\patternSevenInline}{
    % KI: C.10 hält die Inline-Bezeichnung konsistent zum neuen Pattern-Namen.
    \patternnameinline{Structured Audio Companion}
}

%\setlength\parindent{0pt}
%\usepackage{parskip}
\sloppy

\usepackage[nohyperlinks, printonlyused]{acronym}
% Acronym definitions used throughout the paper.
% Keep these definitions before the first \ac{...} usage.

\acrodef{ai}[AI]{artificial intelligence}
\acrodef{adhd}[ADHD]{attention-deficit/hyperactivity disorder}
\acrodef{csv}[CSV]{comma-separated values}
\acrodef{esa}[ESA]{European Space Agency}
\acrodef{gpt}[GPT]{generative pretrained transformer}
\acrodef{llm}[LLM]{large language model}
\acrodef{lms}[LMS]{learning management system}
\acrodef{mllm}[MLLM]{multimodal large language model}
\acrodef{nasa}[NASA]{National Aeronautics and Space Administration}
\acrodef{pdf}[PDF]{portable document format}
\acrodef{svg}[SVG]{Scalable Vector Graphics}
\acrodef{wai}[WAI]{Web Accessibility Initiative}
\acrodef{wcag}[WCAG]{Web Content Accessibility Guidelines}
\acrodef{bvi}[BVI]{blind and visually impaired}
\acrodef{gai}[GenAI]{generative artificial intelligence}
\acrodef{mds}[MDS]{Multimodal Diagrammatic System}
\acrodef{html}[HTML]{Hypertext Markup Language}
\acrodef{kv}[KV]{Karnaugh--Veitch}
\acrodef{mit}[MIT]{Massachusetts Institute of Technology}
\acrodef{oop}[OOP]{Object-Oriented Programming}
\acrodef{spice}[SPICE]{Simulation Program with Integrated Circuit Emphasis}
\acrodef{stem}[STEM]{Science, Technology, Engineering, and Mathematics}
\acrodef{tml}[TML]{textual modeling language}
\acrodef{uml}[UML]{Unified Modeling Language}
\acrodef{uncrpd}[UN-CRPD]{United Nations Convention on the Rights of Persons with Disabilities}
\acrodef{who}[WHO]{World Health Organization}
\acrodef{wysiwyg}[WYSIWYG]{What You See Is What You Get}
\acrodef{xml}[XML]{Extensible Markup Language}
\acrodef{xmi}[XMI]{XML Metadata Interchange}
\acrodef{yaml}[YAML]{Yet Another Markup Language}


\usepackage{xcolor}

\newenvironment{sol}
    {
    %\begin{itshape}
    \begin{bfseries}
    }
    { 
    \end{bfseries}
    \newline
    %\end{itshape}
    }
%

\usepackage[framemethod=TikZ]{mdframed}

\newenvironment{quoted}
    {
    %\begin{itshape}
    }
    { 
    %\end{itshape}
    }
%

\surroundwithmdframed[
    innertopmargin=4pt,
    innerbottommargin=4pt,
    innerrightmargin=4pt,
    innerleftmargin=4pt
]{quoted}

\usepackage{hyperref}
\hypersetup{
    colorlinks=true,
    linkcolor=blue,
    filecolor=magenta,      
    urlcolor=cyan,
    pdftitle={A Pattern Collection for Generating Accessible Teaching Materials in STEM},
    pdfpagemode=FullScreen,
}

% Auswahl der included sections, geht nicht, weil Seitenzahl dann springt
%\includeonly{title_abstract, introduction, 6_Pattern, 7_Pattern, Conclusion_and_Future_Work}
